\documentclass[a4paper]{report}

\makeatletter\@addtoreset{chapter}{part}\makeatother%

\usepackage[margin=2cm]{geometry}
\usepackage{parskip}
\usepackage{amsmath}
\usepackage{amsfonts}

\begin{document}

\begin{titlepage}
	\raggedleft
	\vspace*{2cm}
	\Huge{Mathematics Notes}\\
	\vspace{0.5cm}
	\large{OCR H240}
\end{titlepage}

\tableofcontents

\part{Pure Mathematics}

\chapter{Proof}

\section{Mathematical structures and notation}

Basic mathematical structures:

\begin{itemize}
	\item{Statement/expression: most basic structure. Example: \( x + 3 \).}
	\item{Equation: statements involving equals sign (\( = \)). Only true for discrete values
		of unknown(s). Example: \( x^2 + 3x + 2 = 0 \).}
	\item{Inequality: statements involving inequality sign (\( <, >, \leq, \geq \)). Only true
		for a range of values of the unknown(s). Example: \( 3x + 7 > 4 \).}
	\item{Identity/congruence: statements involving an equivalence sign (\( \equiv, \iff \)).
		True for all values of unknown variable(s). Example:
		\( x^3 - x \equiv x(x + 1)(x - 1) \).}
\end{itemize}

Implication:

\begin{itemize}
	\item{\( A \implies B \): if \( A \) is true then \( B \) is true/\( A \) is sufficient
		for \( B \).}
	\item{\( A \impliedby B \): if \( B \) is true then \( A \) is true/\( A \) is necessary
		for \( B \).}
\end{itemize}

Important sets:

\begin{itemize}
	\item{\( \emptyset \): null set.}
	\item{\( \mathbb{N} \): set of natural numbers (\( 1, 2, 3, \dots \)).}
	\item{\( \mathbb{Z} \): set of integers (\(0, \pm1, \pm2, \pm3, \dots \)).}
	\item{\( \mathbb{R} \):  set of real numbers.} 
\end{itemize}

Set notation:

\begin{itemize}
	\item{\( x > 9 \) becomes \( \{ x : x > 9 \} \).}
	\item{\( x \leq 11 \) becomes \( \{ x : x \leq 11 \} \).}
	\item{\( 3 < x \leq 5 \) becomes \( \{ x : x > 3 \} \cap \{ x : x \leq 5 \} \).}
	\item{\( x > 7 \) or \( x < 2 \) becomes \( \{ x : x > 7 \} \cup \{ x : x < 2 \} \).}
\end{itemize}

Interval notation:

\begin{itemize}
	\item{\( x > 9 \) becomes \( x \in (9, \infty) \).}
	\item{\( x \leq 11 \) becomes \( x \in (-\infty, 11] \).}
	\item{\( 3 < x \leq 5 \) becomes \( x \in (3, 5] \).}
	\item{\( x > 7 \) or \( x < 2 \) becomes \( x \in (7, \infty) \cup (-\infty, 2) \).}
\end{itemize}

Set and interval notation used to show solutions for inequalities. Assumed that
\( x \in \mathbb{R} \) unless otherwise.

\chapter{Algebra and Functions}

\chapter{Coordinate Geometry in the \( x - y \) plane}

\chapter{Sequences and Series}

\chapter{Trigonometry}

\chapter{Exponentials and Logarithms}

\chapter{Differentiation}

\chapter{Integration}

\chapter{Numerical Methods}

\chapter{Vectors}

\part{Statistics}

\chapter{Statistical Sampling}

\chapter{Data Representation and Interpretation}

\chapter{Probability}

\chapter{Statistical Distributions}

\chapter{Statistical Hypothesis Testing}

\part{Mechanics}

\chapter{Quantities and Units in Mechanics}

\chapter{Kinematics}

\chapter{Forces and Newton's Laws}

\chapter{Moments}

\end{document}
